\section{Infinite expectations and triangular arrays}

After treating the ``classical'' case of an i.i.d.
sequence with finite expectations, let's turn to
slightly more exotic situations. First, what happens
if the expectation is infinite? For the strong law,
the following result shows that we have no hope of
having convergence in a meaningful sense.

\begin{thm}[Converse to SLLN]
If $X_1,X_2,\ldots,$ is an i.i.d. sequence of r.v.'s
with $\mathbf{E}|X_1|=\infty$, then

\[
\mathrm{P}\left( \exists \lim_{n\to\infty} \frac{1}{n}\sum_{k=1}^n X_k \right) = 0.
\]

\end{thm}

\begin{proof}
See homework.
\end{proof}

\noindent What about the weak law? It turns out that a weak law
can still hold in certain situations with infinite
expectations, although one has to consider a
generalized empirical average where the sum of the
first $n$ samples is divided by a quantity growing
faster than $n$. Rather than develop a complete
theory, we will consider the particular example of the
$\textbf{St. Petersburg Lottery}$. In this example,
the winning in a single round of lottery is an
integer-valued random variable with the following
distribution:

\[
\mathrm{P}(X=2^k) = 2^{-k}, \qquad k=1,2,3,\ldots
\]

\noindent Let $X_1,X_2,\ldots,$ be an i.i.d. sequence with the
same distribution, and let $S_n = \sum_{k=1}^n X_k$.
How much should you agree to pay to be allowed to play
this lottery $n$ times? (Even the seemingly simple
case $n=1$ of this question  has been a subject of
quite some debate by economists! See \href{http://en.wikipedia.org/wiki/St._Petersburg_paradox}{the Wikipedia
article on the St. Petersburg
paradox}).
In other words, how big should we expect $S_n$ to be,
\emph{with probability close to 1}?

\begin{thm}

\[
\frac{S_n}{n\log_2 n} \xrightarrow[n\to\infty]{\mathrm{P}} 1.
\]

\end{thm}

\noindent In other words, to be allowed to pay the game $n$
times when $n$ is large (and assuming the payoff is
in dollars), it may be considered reasonable to pay
exactly (or even better, slightly less than)
$\log_2 n$ dollars \emph{per round played}. For example,
if $n=1024$ you would be paying \$10 per round.

\begin{proof}
The proof uses a truncation idea similar to the one
we saw in the proof of SLLN, except that for each $n$
we will truncate the first $n$ variables at a level
which is a function of $n$. Denote $b_n = n\log_2 n$,
$Y_{n,k} = X_k \mathbf{1}_{\{X_k<b_n\}}$,
$T_n=\sum_{k=1}^n Y_{n,k}$, and
$a_n = \mathbf{E}(T_n)$. We will prove that

\begin{align*}
\tag{10}\label{eq:thenprove}
\frac{S_n-a_n}{b_n}\xrightarrow[n\to\infty]{\mathrm{P}} 0.
\end{align*}

\noindent First we check that this is enough, by estimating
$a_n$:

\begin{eqnarray*}
a_n &=& \sum_{k=1}^n \mathbf{E}(Y_{n,k}) = \sum_{k=1}^n \mathbf{E}(X_k \mathbf{1}_{\{X_k<b_n\}})
= \sum_{k=1}^n \left(\frac12\cdot 2 + \frac14\cdot 4 + \frac18 \cdot 8 + \ldots + \frac{1}{2^{m_{n}}}\cdot 2^{m_{n}}
\right)
\end{eqnarray*}

\noindent where $m_n$ is the largest integer such that
$2^{m_n} \le b_n$, or in other words
$m_n = \lfloor \log_2 n + \log_2 \log_2 n \rfloor$,
which gives

\[
a_n = \sum_{k=1}^n \lfloor \log_2 n + \log_2 \log_2 n \rfloor = n \log_2 n + O(n \log_2 \log_2 n),
\]

\noindent so $a_n$ indeed behaves like $n\log_2 n$ up to
first-order asymptotics. Now to prove
\eqref{eq:thenprove}, note that for any $\epsilon>0$,

\begin{align*}
\tag{11}\label{eq:lhs}
\mathrm{P}\left( \left|\frac{S_n-a_n}{b_n}\right| > \epsilon\right) \le
\mathrm{P}(T_n\neq S_n) + \mathrm{P}\left( \left|\frac{T_n-a_n}{b_n}\right| > \epsilon\right).
\end{align*}

\noindent In this bound, the first term is bounded by

\[
\sum_{k=1}^n \mathrm{P}(Y_{n,k}\neq X_k) = \sum_{k=1}^n \mathrm{P}(X_k>b_n) \le \sum_{k=1}^n \frac{2}{b_n}
\le \frac{2}{\log_2 n} \xrightarrow[n\to\infty]{} 0.
\]

\noindent To bound the second term, use Chebyshev's inequality
and the fact that
$\mathbf{V}(Y_{n,k})\le \mathbf{E}(Y_{n,k}^2)$ to
write

\begin{eqnarray*}
\mathrm{P}\left( \left|\frac{T_n-a_n}{b_n}\right| > \epsilon\right) &\le& \frac{\mathbf{V}(T_n)}{\epsilon^2 b_n^2}
\le \frac{1}{\epsilon^2 b_n^2} \sum_{k=1}^n \mathbf{E}(Y_{n,k}^2) \\
&\le& \frac{1}{\epsilon^2 b_n^2} \sum_{k=1}^n
\left(\frac12\cdot 2^2 + \frac14\cdot 4^2 + \frac18 \cdot 8^2 + \ldots + \frac{1}{2^{m_{n}}}\cdot 2^{2m_{n}} \right) \\
&\le&
\frac{1}{\epsilon^2 b_n^2} \sum_{k=1}^n 2\cdot 2^{m_n}
\le \frac{2}{\epsilon^2 b_n^2} \sum_{k=1}^n b_n = \frac{2}{\epsilon^2 \log_2 n} \xrightarrow[n\to\infty]{} 0.
\end{eqnarray*}

\noindent We conclude that the left-hand side of \eqref{eq:lhs}
converges to 0, as $n\to\infty$, which finishes the
proof.
\end{proof}

\noindent Another twist on laws of large numbers comes when we
replace the notion of an i.i.d. sequence by a more
general notion of a $\textbf{triangular array}$. In
this case, for each $n$ we have a sequence
$X_{n,1},X_{n,2},\ldots,X_{n,n}$ of independent, but
not necessarily identically distributed, and we denote
$S_n=\sum_{k=1}^n X_{n,k}$ -- this is the sum of the
samples in the $n$-th experiment. Here, for each $n$
there could be a separate experiment involving $n$
different r.v.'s, and the r.v.'s $X_{n,k}$ and
$X_{m,j}$ for $n\neq m$ are not even assumed to be
defined on the same probability space, let alone to be
independent of each other.

Again, instead of giving general conditions for a law
of large numbers to hold, consider the following
example of the so-called
$\textbf{coupon collector's problem}$: A brand of
breakfast cereals comes with a small toy chosen
uniformly at random from a set of $n$ possible kinds
of toys. A collector will buy more boxes of cereals
until she has collected all $n$ different toys.
Denote by $T_n$ the number of boxes she ends up
buying. What can we say about the size of $T_n$?
Fortunately, we can represent it as a sum of
independent r.v.'s, by writing

\[
T_n = X_{n,1} + X_{n,2} + X_{n,3} + \ldots + X_{n,n},
\]

\noindent where

\begin{eqnarray*}
X_{n,1} &=& \textrm{number of boxes purchased to get one kind of toy} \ \ =\ \ 1, \\
X_{n,2} &=& \textrm{number of boxes purchased after having one toy to get a different kind}, \\
X_{n,3} &=& \textrm{number of boxes purchased after having two kinds of toys to get a third kind}, \\
& \vdots&  \\
X_{n,n} &=& \textrm{number of boxes purchased after having $n-1$ kinds of toys to get the last kind}.
\end{eqnarray*}

\noindent Clearly these r.v.'s are independent. Furthermore,
$X_{n,k}$ is a geometric r.v. with parameter
$p_{n,k} = (n-k+1)/n$. This gives us that

\[
\mathbf{E}(T_n) = \sum_{k=1}^n \mathbf{E}(X_{n,k}) = \sum_{k=1}^n \frac{n}{n-k+1} = n\left(\frac{1}{n}+\frac{1}{n-1}+\ldots+\frac{1}{2}+\frac{1}{1}\right) = n H_n,
\]

\noindent where $H_n = \sum_{k=1}^n 1/k$ is the \emph{$n$-th
harmonic number}, and

\[
\mathbf{V}(T_n) = \sum_{k=1}^n \mathbf{V}(X_{n,k}) =
\sum_{k=1}^n \frac{k-1}{n}\left(\frac{n}{n-k+1}\right)^2 \le n^2 \sum_{k=1}^n \frac{1}{k^2} \le 2 n^2
\]

\noindent (in this example, we only need a bound for
$\mathbf{V}(T_n)$, but it is possible also to get
more precise asymptotics for this quantity). It
follows using Chebyshev's inequality that for each
$\epsilon>0$,

\[
\mathrm{P}\left( \left|\frac{T_n-n H_n}{n \log n} \right| > \epsilon\right) \le \frac{\mathbf{V}(T_n)}{\epsilon^2 n^2 (\log n)^2}
\le \frac{2}{\epsilon^2 (\log n)^2} \xrightarrow[n\to\infty]{} 0,
\]

\noindent so $(T_n-n H_n)/(n\log n)$ converges in probability
to $0$, and therefore we get that

\[
\frac{T_n}{n\log n} \xrightarrow[n\to\infty]{\mathrm{P}} 1.
\]

\noindent The lesson to be learned from the above examples is
that some naturally-occurring problems in real life
lead to more complicated situations than can be
modeled with an i.i.d. sequence with finite mean; but
often such problems can be analyzed anyway using the
same ideas and techniques that we developed. In
probability textbooks you can find a general treatment
of various conditions under which a triangular array
of independent random variables satisfies a (strong or
weak) law of large numbers.
